\documentclass[12pt]{article}
\usepackage[margin=1in]{geometry}
\usepackage[T1]{fontenc}
\usepackage{xcolor}
\definecolor{garnet}{HTML}{73000A}

\setlength{\parindent}{0pt}
\setlength{\parskip}{0.2em}

\begin{document}
\noindent
Department of Mechanical Engineering\\
University of South Carolina\\
Columbia, SC 29208\\
\today

\noindent
{\color{garnet}\rule{\textwidth}{0.5pt}}

\vspace{0.2em}

\noindent
Dear Members of the Evaluation Committee:

I recommend Mateo Altomare for the SC Space Grant Undergraduate Research Award. As his faculty advisor in the Integrated Multiphysics and Systems Engineering Laboratory (iMSEL) since June 2025, I have supervised his research on multi-spectral sensor fusion for autonomous maritime perception.

Mateo joined iMSEL as a rising sophomore with no prior experience in machine learning or computer vision. He progressed rapidly from introductory Python to implementing camera calibration routines for our thermal-visible sensor rig, and was independently writing PyTorch training scripts for our Detection Transformer (DETR) model on maritime imagery within his first semester.

He is now co-author on a paper planned for submission at ASME IDETC 2026 presenting a method for pixel-level infrared-to-visible image registration, having derived and implemented the full registration pipeline end-to-end.

Three accomplishments illustrate his technical range. First, he deployed real-time object detection on a power-constrained Raspberry Pi for automated target identification in marine radar displays. Second, he built our multi-spectral detection pipeline that fuses thermal and visible imagery for cross-modal calibration, which now generates all training data for his proposed framework. Third, he curated and annotated approximately 750 thermal-visible frame pairs spanning 1--75\,m depth range with a structured labeling protocol tracking scene conditions, object depth, and registration quality.

His proposed research -- Learned Depth-Adaptive Registration (LDAR) -- addresses a real limitation in our current workflow, which requires knowing object depth in advance. Mateo's proposal would automate depth estimation, eliminating manual specification and additional hardware. I have reviewed his proposed architecture in detail and find the approach to be technically rigorous, well-matched to the available dataset, and targeting an accuracy threshold that, while ambitious, is well within reach based on our preliminary results.

Beyond research, Mateo co-leads the Gamecock Motorsports Formula SAE chassis team and maintains a 3.91\,GPA. Based on his completion of all Fall 2025 milestones ahead of schedule, I am confident he will execute the proposed plan and deliver the stated outcomes.

I give Mateo my highest recommendation for this award.

\vspace{0.5em}

\noindent
Sincerely,

\vspace{0.8em}

\noindent
Dr.\ Yi Wang\\
Associate Professor\\
Department of Mechanical Engineering\\
University of South Carolina\\
\textit{ywang@sc.edu}

\end{document}
